\documentclass{article}
\usepackage{hyperref}
\usepackage{Style}

\nocite{*} % Comentar si quiero citar
%\addbibresource{bibliografia.bib} % Quitar el comentado si quiero usar bibliografia

\begin{document}

\begin{minipage}{2.5cm}
    \includegraphics[width=2cm]{imagen_puc.jpg}
\end{minipage}
\begin{minipage}{12cm}
    {\sc Pontificia Universidad Católica de Chile\\
    Facultad de Matemáticas\\
    Departamento de Matemática\\
    Profesor: Manuel Cabezas -- Ayudante: Benjamín Mateluna}
\end{minipage}
\vspace{2ex}

{\centerline{\bf Introducción a la Geometría - MAT1304}
\centerline{\bf Ayudantía 24}}
\centerline{\bf 16 de junio de 2026}

\vspace{1cm}
\noindent\textbf{Problema 1.} Encuentre el conjunto solución de la ecuación 
$5x^{2}+4xy+2y^{2}-24=0$.

\vspace{5mm}
\noindent\textbf{Problema 2.} Determine que cónica representa cada una de las siguientes 
ecuaciones y obtenga la mayor información posible:

\begin{enumerate}
    \item $4x^{2}+48y+12x=159$
    \item $x^{2}+4y^{2}-10x-40y+109=0$
    \item $9x^{2}-4y^{2}+54x+16y+29=0$
\end{enumerate}

\vspace{1cm}
\noindent\textbf{\underline{Ejercicios Propuestos:}}

\vspace{5mm}
\noindent\textbf{Extra 1.} Demuestre que la cónica de ecuación $xy=1$ es una hipérbole.

\vspace{5mm}
\noindent\textbf{Extra 2.} Determine la ecuación paramétrica de la recta que pasa por los puntos
$(1,0)$ y $(2,\frac{1}{2})$.

%\printbibliography % Quitar el comentado si quiero usar bibliografia

\end{document}